% =========================================================
% Nederpelt & Geuvers, Type Theory and Formal Proof (2014)
% Section 1.4 -- Free and bound variables
% Notes by William Sollers, typeset in the AIM node format.
%
% Page numbers: printed page = PDF page - 29.
% Slots marked \pending{...} are yours to write.
% =========================================================
\documentclass[11pt]{article}
\usepackage{../aim-notes}
\newcommand{\Lam}{\Lambda}

\title{Nederpelt \& Geuvers \S1.4: Free and bound variables}
\author{William Sollers}
\date{Notes, draft 1 (in progress)}

\begin{document}
\maketitle

\begin{remark}[Three kinds of occurrence]\label{rem:occurrences}
\source{N\&G \S1.4 (\textbf{verify} page).}
Variable occurrences in $\lambda$-terms divide into three categories:
\begin{enumerate}[label=(\roman*)]
  \item free occurrences;
  \item bound occurrences;
  \item binding occurrences.
\end{enumerate}
\end{remark}

\begin{example}
\pending{One term with all three kinds, each occurrence labelled. Then the
same variable occurring both free and bound.}
\end{example}

\begin{definition}[Free variables]\label{def:fv}
\pending{N\&G's $\mathrm{FV}$, by recursion on $\Lam$, with its source line.
Compare with $\mathrm{Sub}$ (Def.~1.3.5): same three cases, what differs in
the abstraction case?}
\end{definition}

\end{document}
